\begingroup
\small\setstretch{1.12}\setlength{\tabcolsep}{3pt}
\setlength{\LTleft}{\fill}\setlength{\LTright}{\fill}
\Needspace{13\baselineskip}
\begin{longtable}{@{}>{\raggedright\arraybackslash}p{0.15\linewidth}>{\raggedright\arraybackslash}p{0.15\linewidth}>{\raggedright\arraybackslash}p{0.21\linewidth}>{\raggedright\arraybackslash}p{0.21\linewidth}>{\raggedright\arraybackslash}p{0.21\linewidth}@{}}
\caption{事件標籤下不同靜態成員數之候選平均結果}\label{tab:results-static-size}\\
\toprule
模型數 & 候選數 & AUC & F1-score & Recall \\
\midrule\endfirsthead
\multicolumn{5}{c}{表~\thetable（續）}\\
\toprule
模型數 & 候選數 & AUC & F1-score & Recall \\
\midrule\endhead
\midrule\multicolumn{5}{r}{續下頁}\\\endfoot
\bottomrule\endlastfoot
2 & 9 & 0.7846 & \textbf{0.1071} & 0.1898 \\
3 & 18 & 0.7850 & 0.0973 & \textbf{0.2265} \\
4 & 15 & 0.7859 & 0.0964 & 0.2222 \\
5 & 6 & 0.7866 & 0.1053 & 0.1763 \\
6 & 1 & \textbf{0.7881} & 0.1041 & 0.1941 \\
\end{longtable}
\noindent{\footnotesize (資料來源：本研究事件標籤集成實驗結果)\par 每個種子先平均同一成員數全部候選的 Test 分數，再取十個種子平均。粗體僅供比較，不依此選成員數，也不表示顯著差異。}
\endgroup
